\scanpage{79}
\mspage{39}
l'hom ``local''
\[
(12)\quad \kappa_{i_0}:\quad
\diagram{\begin{tikzcd}[ampersand replacement=\&,row sep=tiny]
T \arrow[r] \arrow[d,equal] \& \pi_1(U_\eta) \\
T_\infty(k) \&
\end{tikzcd}}
\]
associé à l'indice $i_0$. Je vais admettre ce point, qui ne peut guère être difficile.

\margin{\deleted{* En fait, on n'a pas $\unsure{\overline m|\overset{m}{=}}$ besoin que le centre de $\pi$ soit trivial pour cette construction --- c'est un fait général pour les situations fibrées, on trouve une ext. du $\pi_1$ de la base par le $\pi_1$ de la fibre, \unsure{comme}}}

Pour terminer ce n°, je veux encore étudier, dans la situation d'une $U$ \deleted{élémentaire}\inserted{courbe} relative sur un $S$ avec \deleted{\unsure{du}} $U=X\setminus T$, \inserted{$X$ lisse et propre sur $S$,} $T$ fini étale, avec sections $\unsure{g|q}_i$ données de $T$ sur $S$, les ``sections de décomposition'' de l'extension
\[
1-\pi-E_U\longrightarrow E_S-1
\]
associées à $i=i_0$ --- qui définit une \inserted{classe de $\pi$-conj. de} ss-groupe \unsure{fermés|libre} $L_{i_0}$ dans $\pi$. (On suppose qu'on a bien une telle suite exacte, i.e. que $\pi_2(S)\longrightarrow\pi_1(\text{fibre})$ est nul, ce qui sera le cas si $\pi_2(S)=0$, p.ex. (je présume))

\margin{en tous cas, $\unsure{\overline m|\overset{m}{=}}$ sans $\pi_{\unsure{1|2}}(S)-\pi$ trivial, $E_S=\pi_1(S)$ opère extérieurement sur $\pi=\pi_1(\text{fibre géom})$, d'où une extension puisque $\pi$ est à centre trivial. Cette extension admet \unsure{alors} * $\pi_1(U)$ comme quotient\dots}
\margin{il faudrait quand même les vérifier !}

Si on est dans le cas d'un modèle élémentaire au dessus d'un corps de car. zéro (la restriction au car. nulle étant sans doute inutile, si on se borne
\margin{78}
